\documentclass[12pt,a4paper]{amsart}

\usepackage{amsmath,amssymb,amsthm}
\usepackage{mathtools}
\usepackage[utf8]{inputenc}
\usepackage[T1]{fontenc}
\usepackage{hyperref}
\usepackage{enumitem,booktabs,array}
\usepackage{geometry}
\geometry{margin=2.5cm}

% -------------------------------------------------------
% Theorem environments
% -------------------------------------------------------
\newtheorem{theorem}{Theorem}[section]
\newtheorem{lemma}[theorem]{Lemma}
\newtheorem{corollary}[theorem]{Corollary}
\newtheorem{proposition}[theorem]{Proposition}
\theoremstyle{definition}
\newtheorem{definition}[theorem]{Definition}
\newtheorem{example}[theorem]{Example}
\newtheorem{remark}[theorem]{Remark}

% -------------------------------------------------------
% Custom commands
% -------------------------------------------------------
\newcommand{\N}{\mathbb{N}}
\newcommand{\Z}{\mathbb{Z}}
\newcommand{\Q}{\mathbb{Q}}
\newcommand{\R}{\mathbb{R}}
\newcommand{\C}{\mathbb{C}}
\newcommand{\F}{\mathbb{F}}
\newcommand{\A}{\mathbb{A}}
\newcommand{\OK}{\mathcal{O}_K}
\newcommand{\p}{\mathfrak{p}}
\newcommand{\q}{\mathfrak{q}}
\newcommand{\aaa}{\mathfrak{a}}
\newcommand{\bbb}{\mathfrak{b}}
\newcommand{\Nm}{\mathrm{N}}
\newcommand{\disc}{\mathrm{disc}}
\DeclareMathOperator{\Gal}{Gal}
\DeclareMathOperator{\Cl}{Cl}
\DeclareMathOperator{\ord}{ord}
\DeclareMathOperator{\tr}{tr}
\DeclareMathOperator{\nm}{N}

% -------------------------------------------------------
\begin{document}
% -------------------------------------------------------

\title{Algebraic Number Theory:\\
       Rings of Integers, Ideals, the Unit Theorem, and the Class Group}

\author{Michael Fuhrmann}
\address{Specialist Mathematics Project, Build 119}
\email{specialist-maths@localhost}
\date{\today}

\begin{abstract}
We develop the foundations of algebraic number theory, beginning with
number fields and their rings of integers, proceeding through Dedekind's
theory of ideal factorisation, the Minkowski embedding and geometry of
numbers, Dirichlet's Unit Theorem (with complete proof), the class group
and Minkowski's class number bound, and the Dirichlet class number formula.
Special attention is given to quadratic fields $\Q(\sqrt{d})$, where all
concepts can be made fully explicit.  No prior knowledge beyond undergraduate
algebra is assumed.
\end{abstract}

\maketitle
\tableofcontents

% ================================================================
\section{Number Fields and Algebraic Integers}
% ================================================================

\begin{definition}[Number field]
\label{def:number_field}
A \emph{number field} is a finite extension $K/\Q$.  Its \emph{degree}
is $[K:\Q] = n$.
\end{definition}

\begin{definition}[Algebraic integer]
\label{def:alg_int}
An element $\alpha \in \C$ is an \emph{algebraic integer} if it satisfies
a monic polynomial $f \in \Z[x]$:
\[
  \alpha^n + a_{n-1}\alpha^{n-1} + \cdots + a_1\alpha + a_0 = 0,
  \quad a_i \in \Z.
\]
The set of algebraic integers in $K$ is denoted $\OK$.
\end{definition}

\begin{theorem}[Ring of integers]
\label{thm:OK_ring}
$\OK$ is a subring of $K$ with $\OK \cap \Q = \Z$ and $\OK \otimes_\Z \Q = K$.
Moreover, $\OK$ is a free $\Z$-module of rank $n = [K:\Q]$:
\[
  \OK = \Z\omega_1 \oplus \Z\omega_2 \oplus \cdots \oplus \Z\omega_n
\]
for some $\Q$-basis $\{\omega_1, \ldots, \omega_n\}$ of $K$, called an
\emph{integral basis}.
\end{theorem}

\begin{proof}
The fact that algebraic integers are closed under addition and multiplication
follows from linear algebra: if $\alpha$ satisfies a monic polynomial of
degree $m$ and $\beta$ of degree $n$, then $\alpha + \beta$ and $\alpha\beta$
satisfy monic polynomials of degree at most $mn$ (they are eigenvalues of
integer matrices acting on the $\Z$-module $\Z[\alpha,\beta]$, which is
free of rank $mn$).

To see $\OK$ is a free $\Z$-module of rank $n$: any integral basis of $\OK$
is a $\Q$-basis of $K$ (by the rank observation above), so the rank equals
$n$.  The finitely generated $\Z$-module $\OK$ is torsion-free (embedded
in the field $K$), hence free by the classification of finitely generated
abelian groups.
\end{proof}

\begin{example}[Quadratic fields]
\label{ex:quadratic}
Let $d \in \Z$ be squarefree, $d \ne 0, 1$.  Then $K = \Q(\sqrt{d})$
has degree $2$ and:
\[
  \OK =
  \begin{cases}
    \Z\bigl[\tfrac{1+\sqrt{d}}{2}\bigr] = \Z + \Z\tfrac{1+\sqrt{d}}{2}
    & \text{if } d \equiv 1 \pmod{4}, \\[4pt]
    \Z[\sqrt{d}] = \Z + \Z\sqrt{d}
    & \text{if } d \equiv 2, 3 \pmod{4}.
  \end{cases}
\]
\end{example}

\begin{definition}[Discriminant]
\label{def:disc}
Let $\{\omega_1, \ldots, \omega_n\}$ be an integral basis of $\OK$.
The \emph{discriminant} of $K$ (or of $\OK$) is
\[
  \disc(K) = \det\bigl[(\sigma_i(\omega_j))_{i,j}\bigr]^2,
\]
where $\sigma_1, \ldots, \sigma_n : K \hookrightarrow \C$ are the $n$
distinct embeddings of $K$.
\end{definition}

\begin{example}[Discriminant of $\Q(\sqrt{d})$]
For squarefree $d$:
\[
  \disc(\Q(\sqrt{d})) =
  \begin{cases}
    d & \text{if } d \equiv 1 \pmod{4}, \\
    4d & \text{if } d \equiv 2, 3 \pmod{4}.
  \end{cases}
\]
\end{example}

% ================================================================
\section{Dedekind Rings and Ideal Theory}
% ================================================================

\begin{definition}[Dedekind domain]
\label{def:dedekind}
An integral domain $R$ is a \emph{Dedekind domain} if:
\begin{enumerate}[label=(\roman*)]
  \item $R$ is Noetherian,
  \item $R$ is integrally closed in its field of fractions,
  \item every non-zero prime ideal of $R$ is maximal.
\end{enumerate}
\end{definition}

\begin{theorem}
\label{thm:OK_dedekind}
For any number field $K$, the ring $\OK$ is a Dedekind domain.
\end{theorem}

\begin{proof}
(i) \emph{Noetherian:} $\OK$ is a finitely generated $\Z$-module, hence
Noetherian as a $\Z$-module, hence Noetherian as a ring.

(ii) \emph{Integrally closed:} If $\alpha \in K$ satisfies a monic polynomial
over $\OK$, one shows $\alpha \in \OK$ by a standard norm argument.

(iii) \emph{Non-zero primes are maximal:} Let $\p \ne 0$ be prime.
Then $\p \cap \Z = p\Z$ for some rational prime $p$ (since $\OK/\p$
is an integral domain, and $\Z \to \OK/\p$ is injective, forcing $\p \cap \Z$
to be prime, hence $= p\Z$).
Thus $\OK/\p$ is a finite integral domain (it is a quotient of $\OK/p\OK$,
which is a finite ring), hence a field, so $\p$ is maximal.
\end{proof}

\begin{theorem}[Unique factorisation of ideals]
\label{thm:ufid}
Every non-zero ideal $\aaa \subseteq \OK$ has a unique factorisation
\[
  \aaa = \p_1^{e_1} \p_2^{e_2} \cdots \p_r^{e_r},
\]
where $\p_1, \ldots, \p_r$ are distinct non-zero prime ideals and $e_i \ge 1$.
\end{theorem}

\begin{proof}[Proof sketch]
The key ingredient is the existence of \emph{fractional ideals} and their
inverses: for a non-zero prime $\p$ in a Dedekind domain, one constructs
$\p^{-1} = \{x \in K : x\p \subseteq \OK\}$ and shows $\p \cdot \p^{-1} = \OK$.
This makes the non-zero fractional ideals into an abelian group (the
\emph{ideal group}).  Existence of the factorisation uses Noetherianness
(ascending chain condition on ideals); uniqueness uses cancellation in
the ideal group.  See \cite{Neukirch1999}, Chapter I, for the complete argument.
\end{proof}

\begin{definition}[Norm of an ideal]
\label{def:ideal_norm}
For a non-zero ideal $\aaa \subseteq \OK$, its \emph{norm} is
\[
  \Nm(\aaa) = [\OK : \aaa] = |\OK/\aaa|.
\]
For a prime ideal $\p$, $\Nm(\p) = p^f$ where $p = \p \cap \Z$ and $f$
is the \emph{inertia degree} (residue field degree $[\OK/\p : \F_p] = f$).
\end{definition}

\begin{theorem}[Splitting of primes]
\label{thm:splitting}
Let $p \in \Z$ be a rational prime.  In $\OK$:
\[
  p\OK = \p_1^{e_1} \cdots \p_g^{e_g},
\]
with $\sum_{i=1}^g e_i f_i = n$ (the \emph{efg-identity}).
The prime $p$ is:
\begin{itemize}
  \item \emph{split} if all $e_i = 1$ and $g = n$,
  \item \emph{inert} if $g = 1$, $e_1 = 1$, $f_1 = n$,
  \item \emph{ramified} if some $e_i > 1$ (equivalently, $p \mid \disc(K)$).
\end{itemize}
\end{theorem}

% ================================================================
\section{The Minkowski Embedding}
% ================================================================

\begin{definition}[Minkowski embedding]
\label{def:mink_emb}
Let $K$ have $r_1$ real embeddings $\sigma_1, \ldots, \sigma_{r_1} : K \to \R$
and $r_2$ pairs of complex embeddings
$\tau_1, \bar\tau_1, \ldots, \tau_{r_2}, \bar\tau_{r_2} : K \to \C$.
Set $n = r_1 + 2r_2$.  The \emph{Minkowski embedding} is
\[
  \iota : K \longrightarrow \R^n, \quad
  \alpha \longmapsto
  \bigl(\sigma_1(\alpha), \ldots, \sigma_{r_1}(\alpha),
        \Re(\tau_1(\alpha)), \Im(\tau_1(\alpha)), \ldots,
        \Re(\tau_{r_2}(\alpha)), \Im(\tau_{r_2}(\alpha))\bigr).
\]
The image $\iota(\OK)$ is a \emph{full lattice} in $\R^n$ (a discrete
subgroup of rank $n$).
\end{definition}

\begin{theorem}[Minkowski's bound]
\label{thm:mink_bound}
Every ideal class in $\Cl(K)$ contains an integral ideal $\aaa$ with
\[
  \Nm(\aaa) \;\le\; M_K
  \;=\; \frac{n!}{n^n} \left(\frac{4}{\pi}\right)^{r_2} \sqrt{|\disc(K)|}.
\]
In particular, $\Cl(K)$ is \emph{finite} (the \emph{class group} has finite
order $h_K$, the \emph{class number}).
\end{theorem}

\begin{proof}[Proof sketch]
One applies Minkowski's theorem on convex bodies (if a symmetric convex
body $B \subseteq \R^n$ has volume $\text{vol}(B) > 2^n \det(\Lambda)$,
it contains a non-zero lattice point of $\Lambda$) to the lattice $\Lambda$
corresponding to a fractional ideal $\mathfrak{c}^{-1}$ and the convex body
$\{x \in \R^n : \sum |x_i| \le t\}$ (or its modification for complex places).
The explicit volume computation yields the bound $M_K$.
Every ideal class contains some integral ideal of norm $\le M_K$, and the
number of integral ideals of bounded norm is finite, so $h_K < \infty$.
\end{proof}

% ================================================================
\section{The Unit Group: Dirichlet's Unit Theorem}
% ================================================================

\begin{definition}[Unit group]
The \emph{unit group} of $\OK$ is
\[
  \OK^\times = \{u \in \OK : u^{-1} \in \OK\} = \{u \in \OK : \Nm_{K/\Q}(u) = \pm 1\}.
\]
\end{definition}

\begin{theorem}[Dirichlet's Unit Theorem, 1846]
\label{thm:dirichlet_unit}
Let $K$ be a number field with $r_1$ real and $r_2$ pairs of complex
embeddings.  Then
\[
  \OK^\times \;\cong\; \mu(K) \times \Z^{r_1 + r_2 - 1},
\]
where $\mu(K)$ is the (finite cyclic) group of roots of unity in $K$ and
$r = r_1 + r_2 - 1$ is the \emph{unit rank}.
\end{theorem}

\begin{proof}
\textbf{Step 1: The logarithm map.}
Define the \emph{Minkowski logarithm}:
\[
  \ell : \OK^\times \longrightarrow \R^{r_1+r_2},
  \quad
  u \mapsto \bigl(\log|\sigma_1(u)|, \ldots, \log|\sigma_{r_1}(u)|,
               2\log|\tau_1(u)|, \ldots, 2\log|\tau_{r_2}(u)|\bigr).
\]
Since $\prod_i |\sigma_i(u)|^{} \cdot \prod_j |\tau_j(u)|^2 = |N_{K/\Q}(u)| = 1$,
the image of $\ell$ lands in the hyperplane
\[
  H = \bigl\{(x_1,\ldots,x_{r_1+r_2}) : x_1 + \cdots + x_{r_1+r_2} = 0\bigr\}
  \cong \R^{r_1+r_2-1}.
\]

\textbf{Step 2: The kernel.}
An element $u \in \OK^\times$ satisfies $\ell(u) = 0$ iff $|\sigma_i(u)| = 1$
for all $i$.  By Minkowski's bound, there are finitely many units of bounded norm,
so $\ker(\ell)$ is finite.  A finite subgroup of $K^\times$ is cyclic (every finite
subgroup of the multiplicative group of a field is cyclic).  Hence $\ker(\ell) = \mu(K)$.

\textbf{Step 3: The image is a full lattice of rank $r_1+r_2-1$.}
We must show $\ell(\OK^\times)$ is a full lattice in $H$.

\emph{Discreteness:}  The set of units $u$ with $\|\ell(u)\|_\infty \le B$
satisfies $e^{-B} \le |\sigma_i(u)| \le e^B$ for all $i$, so all
conjugates are bounded.  By the boundedness of $\OK$ within the Minkowski
space, there are finitely many such units.

\emph{Full rank:}  We must find $r_1+r_2-1$ linearly independent elements
in $\ell(\OK^\times)$.  For each embedding $\sigma_k$, one constructs a unit
with $|\sigma_k(u)| > 1$ and $|\sigma_j(u)| < 1$ for $j \ne k$ using
geometry of numbers (apply Dirichlet's box principle to lattice points in
a suitable box).  The resulting vectors $\ell(u_1), \ldots, \ell(u_{r_1+r_2})$
span $H$ (they satisfy the hyperplane equation), and $r_1+r_2-1$ of them
are linearly independent.

\textbf{Conclusion:}  The exact sequence
$1 \to \mu(K) \to \OK^\times \xrightarrow{\ell} \ell(\OK^\times) \to 0$
with $\ell(\OK^\times) \cong \Z^{r_1+r_2-1}$ (free abelian, being a lattice)
gives $\OK^\times \cong \mu(K) \times \Z^{r_1+r_2-1}$.
\end{proof}

\begin{example}[Quadratic fields]
For $K = \Q(\sqrt{d})$:
\begin{itemize}
  \item Real quadratic ($d > 0$): $r_1 = 2$, $r_2 = 0$, rank $= 1$.
    So $\OK^\times \cong \{\pm 1\} \times \Z = \{\pm \varepsilon^n : n \in \Z\}$
    for a \emph{fundamental unit} $\varepsilon > 1$.
  \item Imaginary quadratic ($d < 0$): $r_1 = 0$, $r_2 = 1$, rank $= 0$.
    So $\OK^\times$ is finite:
    $\OK^\times = \{\pm 1, \pm i\}$ for $d=-1$ (Gaussian integers),
    $\OK^\times = \{\pm 1, \pm\omega, \pm\omega^2\}$ ($\omega = e^{2\pi i/3}$)
    for $d=-3$, and $\OK^\times = \{\pm 1\}$ for all other $d < 0$.
\end{itemize}
\end{example}

\begin{example}[Fundamental unit of $\Q(\sqrt{2})$]
$d=2$, $\OK = \Z[\sqrt{2}]$, $\varepsilon = 1 + \sqrt{2}$.
Check: $(1+\sqrt{2})(1-\sqrt{2}) = 1 - 2 = -1$, so $N(\varepsilon) = -1$
and $\varepsilon^{-1} = -(1-\sqrt{2}) = \sqrt{2}-1$.
\end{example}

% ================================================================
\section{The Class Group}
% ================================================================

\begin{definition}[Fractional ideals and the class group]
A \emph{fractional ideal} of $\OK$ is a non-zero $\OK$-submodule
$\mathfrak{I} \subseteq K$ such that $d\mathfrak{I} \subseteq \OK$
for some $d \in \Z \setminus \{0\}$.  The fractional ideals form an
abelian group $\mathcal{I}_K$ under multiplication.  The \emph{principal}
fractional ideals $\alpha\OK$ ($\alpha \in K^\times$) form a subgroup
$\mathcal{P}_K$.  The \emph{ideal class group} is
\[
  \Cl(K) = \mathcal{I}_K / \mathcal{P}_K.
\]
Its order $h_K = |\Cl(K)|$ is the \emph{class number}.
\end{definition}

\begin{theorem}
$h_K = 1$ if and only if $\OK$ is a principal ideal domain (equivalently,
a unique factorisation domain).
\end{theorem}

\begin{proof}
In a Dedekind domain, every ideal is principal iff the class group is trivial.
A Dedekind domain is a UFD iff it is a PID (since every irreducible element
generates a prime ideal in a Dedekind domain).
\end{proof}

\begin{example}[Class number 1 imaginary quadratic fields]
The imaginary quadratic fields with $h_K = 1$ are exactly those with
$d \in \{-1, -2, -3, -7, -11, -19, -43, -67, -163\}$.
This is the \emph{Stark--Heegner theorem} (1967/1969).
\end{example}

\begin{example}[Class group of $\Q(\sqrt{-5})$]
$\disc = -20$, Minkowski bound $M_K = \frac{2}{\pi}\sqrt{20} \approx 2.85$.
We check primes $p \le 2$: The prime $2$ ramifies: $2\OK = (2, 1+\sqrt{-5})^2$.
Let $\p = (2, 1+\sqrt{-5})$.  Then $\Nm(\p) = 2$, $\p \notin \mathcal{P}_K$
(since $2 = a^2 + 5b^2$ has no integer solution), but $\p^2 = 2\OK$ is principal.
Hence $\Cl(K) \cong \Z/2\Z$ and $h_K = 2$.
\end{example}

% ================================================================
\section{Splitting of Primes in Quadratic Fields}
% ================================================================

\begin{theorem}[Splitting via Legendre symbol]
\label{thm:quad_split}
Let $K = \Q(\sqrt{d})$ with discriminant $D = \disc(K)$.  An odd prime $p$
satisfies:
\[
  p\OK =
  \begin{cases}
    \p\q, \quad \p \ne \q & \text{if } \bigl(\tfrac{D}{p}\bigr) = +1
    \quad\text{(split)}, \\[2pt]
    \p^2 & \text{if } p \mid D
    \quad\text{(ramified)}, \\[2pt]
    \p \text{ (prime)} & \text{if } \bigl(\tfrac{D}{p}\bigr) = -1
    \quad\text{(inert)}.
  \end{cases}
\]
\end{theorem}

\begin{proof}
The minimal polynomial of $\sqrt{d}$ over $\Q$ is $f(x) = x^2 - d$.
Splitting of $p$ is determined by the factorisation of $f \pmod{p}$.
By quadratic reciprocity and the definition of the Legendre symbol,
$f$ splits into two distinct linear factors mod $p$ iff $d$ is a quadratic
residue mod $p$ (iff $\bigl(\tfrac{d}{p}\bigr) = 1$), and is irreducible
mod $p$ iff $d$ is a non-residue.  Since $D$ and $d$ differ only by a
square factor (which does not affect the Legendre symbol for odd primes),
$\bigl(\tfrac{d}{p}\bigr) = \bigl(\tfrac{D}{p}\bigr)$.
\end{proof}

% ================================================================
\section{The Class Number Formula}
% ================================================================

\begin{definition}[Dedekind zeta function]
\label{def:dedekind_zeta}
For a number field $K$, the \emph{Dedekind zeta function} is
\[
  \zeta_K(s) = \sum_{\substack{\aaa \subseteq \OK \\ \aaa \ne 0}} \frac{1}{\Nm(\aaa)^s}
  = \prod_{\p} \frac{1}{1 - \Nm(\p)^{-s}}, \qquad \Re(s) > 1.
\]
It has analytic continuation to $\C \setminus \{1\}$ with a simple pole at $s=1$.
\end{definition}

\begin{theorem}[Dirichlet--Dedekind Class Number Formula]
\label{thm:class_number_formula}
The residue of $\zeta_K(s)$ at $s = 1$ is
\[
  \lim_{s \to 1^+} (s-1)\,\zeta_K(s)
  \;=\;
  \frac{2^{r_1} (2\pi)^{r_2}\, h_K\, R_K}{w_K \sqrt{|\disc(K)|}},
\]
where $h_K$ is the class number, $R_K$ is the \emph{regulator}
($|\det$ of the $r \times r$ submatrix of the logarithm map applied to a
system of fundamental units$|$), and $w_K = |\mu(K)|$ is the number of roots
of unity.
\end{theorem}

\begin{proof}[Proof sketch]
The Dedekind zeta function factors as a product over ideal classes.
For each class $[\mathfrak{c}] \in \Cl(K)$, one sums over ideals in that class:
\[
  \zeta_K(s) = \sum_{[\mathfrak{c}] \in \Cl(K)} \zeta_K(s, [\mathfrak{c}]),
  \qquad
  \zeta_K(s, [\mathfrak{c}]) = \sum_{\substack{\aaa \in [\mathfrak{c}] \\ \aaa \ne 0}}
  \Nm(\aaa)^{-s}.
\]
Each partial zeta function $\zeta_K(s,[\mathfrak{c}])$ has residue
\[
  \frac{2^{r_1}(2\pi)^{r_2} R_K}{w_K \sqrt{|\disc(K)|}}
\]
(computed by counting lattice points in the Minkowski space, using the
Minkowski embedding and the Volume of the fundamental domain), and there
are $h_K$ classes.  Summing over all classes gives the stated formula.
See \cite{Neukirch1999}, Chapter I, \S8, for the complete proof.
\end{proof}

\begin{corollary}[Class number of imaginary quadratic fields]
For $K = \Q(\sqrt{d})$, $d < 0$ squarefree: $r_1 = 0$, $r_2 = 1$, $R_K = 1$,
$\disc(K) < 0$.  The formula gives
\[
  h_K = \frac{w_K \sqrt{|D|}}{2\pi} L(1, \chi_D),
\]
where $\chi_D = \bigl(\tfrac{D}{\cdot}\bigr)$ is the Kronecker symbol and
$L(s, \chi_D) = \sum_{n=1}^\infty \chi_D(n) n^{-s}$ is the Dirichlet $L$-function.
\end{corollary}

% ================================================================
\section{Worked Examples: Class Groups of Small Quadratic Fields}
% ================================================================

\begin{example}[$\Q(\sqrt{-23})$]
$D = -23$ (since $-23 \equiv 1 \pmod{4}$).  Minkowski bound:
$M_K = \frac{2}{\pi}\sqrt{23} \approx 3.05$.
Check $p = 2$: $\bigl(\tfrac{-23}{2}\bigr) = 1$ (since $-23 \equiv 1 \pmod 8$),
so $2$ splits: $2\OK = \p_2 \q_2$.  Norm of $\p_2 = 2$.
Check $p = 3$: $\bigl(\tfrac{-23}{3}\bigr) = \bigl(\tfrac{1}{3}\bigr) = 1$,
so $3$ splits.
One verifies that $\p_2 \notin \mathcal{P}_K$ but $\p_2^3 \in \mathcal{P}_K$
(since $N(a + b\tfrac{1+\sqrt{-23}}{2}) = a^2 + ab + 6b^2$, which represents
$8$ at $a=1, b=1$ but not $2$ or $4$).
Hence $h_{-23} = 3$ and $\Cl(K) \cong \Z/3\Z$.
\end{example}

\begin{example}[Real quadratic $\Q(\sqrt{5})$]
$d = 5 \equiv 1 \pmod{4}$, so $\OK = \Z[\phi]$ with $\phi = \tfrac{1+\sqrt{5}}{2}$
(the golden ratio).  $D = 5$.
Fundamental unit: $\phi = \tfrac{1+\sqrt{5}}{2}$, since $N(\phi) = \phi\bar\phi
= \tfrac{(1+\sqrt{5})(1-\sqrt{5})}{4} = \tfrac{1-5}{4} = -1$.
$h_{\Q(\sqrt{5})} = 1$ (Minkowski bound $< 2$, so trivial class group).
\end{example}

% ================================================================
\section{Summary and Further Reading}
% ================================================================

\begin{remark}[What comes next]
The theory developed here is the starting point for:
\begin{itemize}
  \item \textbf{Iwasawa theory} (Paper 38): $p$-adic $L$-functions arise from
    varying the class group in towers of number fields.
  \item \textbf{Langlands programme}: $L$-functions of Galois representations
    of $\Gal(\bar\Q/\Q)$ generalise the Dedekind zeta function.
  \item \textbf{Class field theory}: The Artin map identifies $\Cl(K)$ with
    the Galois group of the Hilbert class field $H_K/K$.
  \item \textbf{Explicit class field theory} for CM fields is the basis of the
    theory of complex multiplication of elliptic curves.
\end{itemize}
\end{remark}

\begin{thebibliography}{99}

\bibitem{Neukirch1999}
J.~Neukirch,
\textit{Algebraic Number Theory},
Springer, Berlin, 1999.

\bibitem{Marcus1977}
D.~A.~Marcus,
\textit{Number Fields},
Springer, New York, 1977.

\bibitem{Milne2020}
J.~S.~Milne,
\textit{Algebraic Number Theory} (version 3.08),
\url{https://www.jmilne.org/math/CourseNotes/ANT.pdf}, 2020.

\bibitem{Samuel1970}
P.~Samuel,
\textit{Algebraic Theory of Numbers},
Hermann, Paris, 1970.

\bibitem{Cohen1993}
H.~Cohen,
\textit{A Course in Computational Algebraic Number Theory},
Springer, Berlin, 1993.

\bibitem{Washington1997}
L.~C.~Washington,
\textit{Introduction to Cyclotomic Fields},
2nd ed., Springer, New York, 1997.

\end{thebibliography}

\end{document}
